%Sprache
\usepackage[ngerman]{babel}
\setkomafont{disposition}{\normalcolor \bfseries}
\usepackage[autostyle=true,german=quotes]{csquotes}

%Schrift
\usepackage[T1]{fontenc}
\usepackage{lmodern}

%Mathepakete
\usepackage{amsmath}
\usepackage{amssymb}
\usepackage{amsthm} %Mehr Optionen für \newtheorem wir Stile, und proof Umgebung
%\usepackage{siunitx}  
%\sisetup{locale = DE} 
%\sisetup{separate-uncertainty=true}

%Paket für Hyperlinks
\usepackage{hyperref}

%Einbindung von Bildern
\usepackage{graphicx}
\usepackage{float}
\usepackage{wrapfig}
\addto\captionsngerman{\renewcommand{\figurename}{Abb.}}
%richtige Angabe von Quellen
\usepackage[backend=biber, sorting=none]{biblatex} 
\addbibresource{Einstellungen/Quellen.bib}
%Neudefinitionen
\usepackage{layouts} %ermittelt die Dimensionen der aktuellen Seiten
\newcommand{\druckegroesse}{
    Breite: \printinunitsof{cm}\prntlen{\textwidth}\\
    Höhe: \printinunitsof{cm}\prntlen{\textheight}
}

%Mathematik: Sätze, Lemmate, Definitionen, Bemerkungen, Beispiele ...
\newtheoremstyle{nsatz} % Name
  {0pt}             % Space above
  {\topsep}             % Space below
  {\itshape}                    % Body font
  {\parindent}                    % Indent amount
  {\bfseries}           % Head font
  {:}                   % Punctuation after head
  {\newline}            % Space after head  
  {}                    % Head spec
\theoremstyle{nsatz}
\newtheorem{satz}{Satz}[section]
\newtheorem{korollar}{Korollar}[satz]
\newtheorem{lemma}[satz]{Lemma}
\newtheorem{beispiel}[satz]{Beispiel}
\newtheorem{beispiele}[satz]{Beispiele}
\newtheorem{defusatz}[satz]{Definition und Satz}


\newtheoremstyle{ndefinition} % Name
  {0pt}             % Space above
  {\topsep}             % Space below
  {}                    % Body font
  {\parindent}                    % Indent amount
  {\bfseries}           % Head font
  {:}                   % Punctuation after head
  {\newline}            % Space after head  
  {}                    % Head spec
\theoremstyle{ndefinition}
\newtheorem{definition}[satz]{Definition}

\theoremstyle{remark}
\newtheorem*{bemerkung}{Bemerkung}

%Mathematik: Beweise
\renewcommand\qedsymbol{}
\newenvironment{beweis}
  {\begin{proof}}
  {\par\noindent\makebox[\linewidth]{q.e.d.}\end{proof}}

%Vereinfachte Einbindung der Vorlesungen
%Die Vorlesungen müssen in einem Ordner "Vorlesungen" liegen, und müssen "Vorlesung01.tex", "Vorlesung02.tex", ... heißen. Dann können sie mit \vorlesungen{1}{10} eingebunden werden.
\usepackage{pgffor}%
\newcommand{\vorlesungen}[2]{
	\foreach \c in {#1,...,#2}{
		\ifnum\c<10
			\edef\nummer{0\c}
		\else
			\edef\nummer{\c}
		\fi
		\IfFileExists{Vorlesungen/Vorlesung\nummer.tex}{
			\input{Vorlesungen/Vorlesung\nummer.tex}
		}{}%
	}%
}